% Created 2024-01-08 lun 09:17
% Intended LaTeX compiler: pdflatex
\documentclass[11pt]{article}
\usepackage[utf8]{inputenc}
\usepackage[T1]{fontenc}
\usepackage{graphicx}
\usepackage{longtable}
\usepackage{wrapfig}
\usepackage{rotating}
\usepackage[normalem]{ulem}
\usepackage{amsmath}
\usepackage{amssymb}
\usepackage{capt-of}
\usepackage{hyperref}
\usepackage{minted}

\usepackage{parskip}
\usepackage[utf8]{inputenc} %% For unicode chars
\usepackage{placeins}
\usepackage[margin=2.50cm]{geometry}
\usepackage[T1]{fontenc}
\usepackage{mathpazo}
\linespread{1.05}
\usepackage[scaled]{helvet}
\usepackage{courier}
\hypersetup{colorlinks=true,linkcolor=blue}
\RequirePackage{fancyvrb}
\DefineVerbatimEnvironment{verbatim}{Verbatim}{fontsize=\small,formatcom = {\color[rgb]{0.5,0,0}}}
\usepackage{mdframed}
\BeforeBeginEnvironment{minted}{\begin{mdframed}}
\AfterEndEnvironment{minted}{\end{mdframed}}
\setcounter{secnumdepth}{3}
\author{José Manuel Sánchez Álvarez}
\date{}
\title{PHP CRUD Car Dealership. Update.}
\hypersetup{
 pdfauthor={José Manuel Sánchez Álvarez},
 pdftitle={PHP CRUD Car Dealership. Update.},
 pdfkeywords={},
 pdfsubject={},
 pdfcreator={Emacs 28.1 (Org mode 9.5.5)}, 
 pdflang={Spanish}}
\begin{document}

\maketitle
\setcounter{tocdepth}{3}
\tableofcontents

\setlength\parindent{10pt}
To add the functionality of deleting a car from your CRUD application,
you'll need both a PHP script to handle the deletion process and an HTML
interface for users to trigger this action. Below is a basic example of
how you can implement this.

\section{Step 1: HTML Interface for Car Deletion}
\label{sec:org048d4fe}
You can provide a simple form for users to input the ID of the car they
wish to delete. This form will send the data to a PHP script
(\texttt{delete\_car.php}) for processing.

\textbf{HTML (part of your existing interface, e.g., \texttt{delete\_car.html}):}

\begin{minted}[]{php}
<!DOCTYPE html>
<html>
<head>
    <title>Delete Car</title>
</head>
<body>
    <h1>Delete Car</h1>
    <form action="delete_car.php" method="post">
        Car ID to Delete: <input type="text" name="car_id" required>
        <input type="submit" value="Delete Car">
    </form>
</body>
</html>
\end{minted}

\section{Step 2: PHP Script for Deleting a Car}
\label{sec:orgbcfed8d}
This script (\texttt{delete\_car.php}) will receive the car ID from the form and
use it to delete the corresponding record from the database.

\textbf{PHP (delete\_car.php):}

\begin{minted}[]{php}
<?php
require_once 'db_connect.php'; // Database connection file

if ($_SERVER["REQUEST_METHOD"] == "POST" && !empty($_POST['car_id'])) {
    $carId = $_POST['car_id'];

    // Prepare a delete statement
    $stmt = $pdo->prepare("DELETE FROM Car WHERE id = :carId");
    $stmt->bindParam(':carId', $carId, PDO::PARAM_INT);

    try {
        $stmt->execute();
        if ($stmt->rowCount() > 0) {
            echo "Car with ID $carId has been deleted.";
        } else {
            echo "No car found with ID $carId.";
        }
    } catch (PDOException $e) {
        echo "Error deleting car: " . $e->getMessage();
    }
} else {
    echo "Car ID is required.";
}
?>
\end{minted}

\section{Explanation:}
\label{sec:org9137f5d}
\begin{enumerate}
\item \textbf{Form Submission}: The HTML form provides a field for the user to
enter the ID of the car to be deleted.

\item \textbf{PHP Processing}:

\begin{itemize}
\item The PHP script checks if the form has been submitted and if the
\texttt{car\_id} field is not empty.
\item It then prepares a SQL DELETE statement to remove the car from the
database.
\item The car ID from the form is bound to the prepared statement and
executed.
\item The script checks if a row was affected (i.e., a car was deleted)
and provides appropriate feedback.
\end{itemize}
\end{enumerate}

\section{Security and Error Handling:}
\label{sec:orga2955ba}
\begin{itemize}
\item \textbf{Prepared Statements}: This method prevents SQL injection attacks.
\item \textbf{Error Handling}: Try-catch blocks are used to handle any exceptions
that may occur during the database operation.
\item \textbf{Validation}: The script checks if the \texttt{car\_id} field is filled before
attempting a delete operation.
\end{itemize}

\section{Integration:}
\label{sec:org39c5f88}
\begin{itemize}
\item Integrate this functionality into your existing CRUD application,
ensuring that the database connection file (\texttt{db\_connect.php}) and any
other required files are correctly included.
\item Enhance user experience by adding confirmation prompts or messages,
and redirecting or refreshing the page as needed after deletion.
\end{itemize}

This setup provides basic functionality for deleting a car from the
database. For a production-level application, you might want to
implement more robust validation, error handling, and user feedback
mechanisms.
\end{document}
